\section{Plantilla de formalizaci\'on de requisitos}
\label{sec:anexo-plantilla}

La siguiente plantilla se utiliz\'o para formalizar todos los requisitos del corpus.

\begin{lstlisting}[language=,caption={Plantilla de requisitos},label=lst:plantilla]
ID:
REQ-[DOMINIO]-[NUMERO]

Nombre:

Tipo:
[Funcional | Calidad | Restriccion]

Descripcion formal:
El sistema debera [accion] [objeto] [condicion]

Fuente:

Rationale:

Prioridad:
[Alta | Media | Baja]

Verificacion:

Estado:
[Borrador | Elaborado | Validado | Implementado | Archivado]

Version:

Dependencias:

Modulo:
\end{lstlisting}

\subsection*{Reglas obligatorias para la descripci\'on formal}

\begin{enumerate}
    \item Debe expresar una \'unica obligaci\'on.
    \item Debe comenzar por la f\'ormula ``El sistema deber\'a...''.
    \item Debe contener un verbo observable y verificable.
    \item Debe evitar t\'erminos subjetivos o ambiguos.
    \item Debe incluir condiciones cuantificables cuando sea posible.
    \item Debe poder verificarse mediante prueba, an\'alisis, inspecci\'on o demostraci\'on.
    \item Debe ser comprensible sin necesidad de interpretar otros requisitos.
    \item No debe describir detalles de implementaci\'on salvo que se trate de una restricci\'on expl\'icita.
\end{enumerate}
